\documentclass[11pt,a4paper]{article}
\usepackage[T1]{fontenc}
\usepackage[utf8]{inputenc}
\usepackage[brazilian]{babel}
\usepackage{lmodern}
\usepackage[margin=2.2cm]{geometry}
\usepackage{enumitem}
\usepackage{booktabs}
\usepackage{xcolor}
\usepackage{tcolorbox}
\usepackage{fancyhdr}
\usepackage{amssymb}

\pagestyle{fancy}
\fancyhf{}
\lhead{\small Roteiro de entrevista com gestores}
\rhead{\small TCC, César Hoffmann, UNISINOS}
\cfoot{\small\thepage}

\setlength{\parindent}{0pt}
\setlength{\parskip}{0.5em}

\newcommand{\perfil}[1]{{\unskip\nobreak\hfil\penalty50\hskip2em\hbox{}\nobreak\hfill\mbox{\small\textbf{\textcolor{gray}{[#1]}}}\parfillskip=0pt\finalhyphendemerits=0\par}}
\newcommand{\aux}[1]{\par\vspace{0.15em}{\small\textcolor{gray}{\textit{Aprofundar: #1}}}}
\newcommand{\bloco}[3]{\section*{#1 \hfill{\normalsize\textcolor{gray}{#2 \textbullet{} #3}}}}
\newcommand{\caixa}{$\square$\ }

\setlist[enumerate]{label=\textbf{\arabic*.}, leftmargin=*, itemsep=0.8em}

\begin{document}

\begin{center}
  {\Large\bfseries Roteiro de entrevista com gestores}\\[0.4em]
  Cardápios digitais e autoatendimento em bares e restaurantes de Porto Alegre\\[0.2em]
  {\small César Hoffmann \textbullet{} orientação: Prof.\ Dr.\ Guilherme Lacerda \textbullet{} Ciência da Computação, UNISINOS}
\end{center}

\begin{tcolorbox}[colback=gray!6,colframe=gray!50,title=Antes de começar,fonttitle=\bfseries]
\caixa Gravador testado e com bateria (gravar também no celular como reserva).\\
\caixa TCLE impresso em duas vias, assinado antes de ligar o gravador.\\
\caixa Termo de Confidencialidade, se o gestor for falar de custos, contratos ou números.\\
\caixa Termo de Uso de Imagem, se for fotografar o espaço, o totem ou as torneiras.\\
\caixa Perguntar se o nome do estabelecimento pode aparecer no trabalho ou se prefere pseudônimo.\\
\caixa Ao ligar o gravador: dizer data, local, perfil do caso e pedir o ``ok'' verbal de consentimento.
\end{tcolorbox}

\textbf{Duração-alvo:} 45 a 60 minutos. \textbf{Como usar:} a etiqueta à direita de cada pergunta indica para quais perfis ela vale. Perguntas em cinza são aprofundamentos opcionais; use quando a resposta vier curta ou genérica. Sempre que possível, peça um exemplo concreto (``lembra de uma vez em que\ldots?'').

\begin{center}
\small
\begin{tabular}{lll}
  \toprule
  \textbf{Perfil} & \textbf{Caso} & \textbf{Blocos aplicáveis} \\
  \midrule
  A, pré-2020 que adotou & Mambembe Cervejaria & 0, 1, 2, 3, 5, 6, 7 \\
  B, pré-2020 que resistiu & Lancheria do Parque & 0, 1, 4, 5, 6, 7 \\
  C, nascido digital & Cervejaria Fora da Lei (Assis Brasil) & 0, 2, 3, 5, 6, 7 \\
  D, pós-2020 que optou pelo físico & Ordinária Pub (Bom Fim) & 0, 1*, 4, 5, 6, 7 \\
  \bottomrule
\end{tabular}
\end{center}

\textbf{Perfil D.} Use as perguntas marcadas com B. *No Bloco 1, faça só a pergunta 9, trocada por: ``Ao abrir a casa, vocês consideraram usar cardápio digital? Por que escolheram só o físico?''.

\textbf{Termos usados com o entrevistado.} Fale em ``sistema de autoatendimento'' ou ``cardápio digital'', adaptando ao caso: QR Code, tablet, totem ou torneiras de autosserviço com crédito no aplicativo ou cartão.

%-------------------------------------------------------------
\bloco{Bloco 0, Caracterização}{todos}{5 min}
\begin{enumerate}[series=perguntas]
  \item Conte brevemente a história da casa: quando abriu, qual é o conceito e quem é o cliente típico. \perfil{A B C}
  \item Qual é o seu papel aqui e há quanto tempo você está nessa função? \perfil{A B C}
  \item Mais ou menos quantas mesas, quantos funcionários por turno e qual o tíquete médio? \perfil{A B C}
  \item Hoje, como o cliente consulta o cardápio e faz o pedido? Há quanto tempo funciona assim? \perfil{A B C}
\end{enumerate}

%-------------------------------------------------------------
\bloco{Bloco 1, Trajetória}{A e B}{8 min}
\begin{enumerate}[resume=perguntas]
  \item Como funcionavam o cardápio e o pedido antes de 2020? \perfil{A B}
  \item O que mudou, ou não mudou, durante e depois da pandemia? \perfil{A B}
  \item O que motivou a decisão de adotar o autoatendimento? Quem decidiu? Foi algo planejado ou emergencial? \perfil{A}
    \aux{havia pressão de clientes, concorrentes, fornecedores ou da falta de funcionários?}
  \item Como foi a transição? Quanto tempo levou e o que foi mais difícil (equipe, clientes, tecnologia)? \perfil{A}
  \item Por que vocês decidiram manter o atendimento tradicional? Já chegaram a considerar mudar? \perfil{B}
\end{enumerate}

%-------------------------------------------------------------
\bloco{Bloco 2, Solução adotada e arquitetura}{A e C}{12 min}
\begin{enumerate}[resume=perguntas]
  \item Que modalidade vocês usam (QR Code no celular do cliente, tablet, totem, torneiras com crédito) e por que essa? \perfil{A C}
  \item Qual sistema ou fornecedor vocês usam? Por que esse? Já trocaram de fornecedor? \perfil{A C}
  \item Vocês chegaram a pensar em desenvolver um sistema próprio em vez de contratar um pronto? Por quê? \perfil{A C}
  \item O sistema conversa com outros sistemas da casa, como caixa (PDV), estoque, emissão de nota ou pedidos da cozinha? Isso é automático ou alguém precisa redigitar? \perfil{A C}
  \item O que acontece quando a internet ou o sistema cai? Já aconteceu num dia de movimento? Como resolveram? \perfil{A C}
    \aux{existe algum plano B (cardápio impresso, anotação manual, modo offline)?}
  \item Como e por quem o cardápio é atualizado (preços, itens novos, itens esgotados)? Quanto tempo leva? \perfil{A C}
  \item Qual o custo aproximado (mensalidade, taxas por transação, equipamentos) e como você avalia o custo-benefício? \perfil{A C}
\end{enumerate}

%-------------------------------------------------------------
\bloco{Bloco 3, Operação e requisitos}{A e C}{8 min}
\begin{enumerate}[resume=perguntas]
  \item Quais problemas técnicos aparecem com mais frequência (conexão, QR Code que não abre, tela que não responde, cliente que não entende a interface)? \perfil{A C}
  \item Na sua visão, o que um bom sistema desses \emph{precisa} fazer? E o que ele não pode deixar de ter em termos de velocidade, estabilidade ou facilidade? \perfil{A C}
    \aux{peça para listar ``três coisas indispensáveis'' e ``uma coisa que irrita''.}
  \item Como o autoatendimento afetou o tempo de atendimento, o tíquete médio, o tamanho da equipe e os erros de pedido? \perfil{A C}
  \item O sistema pede cadastro ou dados pessoais do cliente (nome, CPF, telefone)? Para quê? Algum cliente já reclamou disso? \perfil{A C}
\end{enumerate}

%-------------------------------------------------------------
\bloco{Bloco 4, Resistência e barreiras}{B}{8 min}
\begin{enumerate}[resume=perguntas]
  \item Que receios você tem em relação ao cardápio digital ou ao autoatendimento (custo, depender de fornecedor, perder o contato com o cliente)? \perfil{B}
  \item Como acha que seus clientes reagiriam? Já fizeram algum teste? \perfil{B}
  \item Já foram procurados por fornecedores desse tipo de sistema? O que ofereceram e por que não fechou? \perfil{B}
  \item O que precisaria acontecer para vocês adotarem? \perfil{B}
\end{enumerate}

%-------------------------------------------------------------
\bloco{Bloco 5, Percepção do cliente e acessibilidade}{todos}{8 min}
\begin{enumerate}[resume=perguntas]
  \item Que comentários, positivos e negativos, vocês recebem dos clientes sobre o cardápio ou o pedido? \perfil{A B C}
  \item Vocês percebem diferença entre tipos de cliente (idade, turistas, grupos grandes, casais)? \perfil{A B C}
  \item Há clientes com dificuldade para usar o sistema, como pessoas idosas ou com baixa visão? Como a equipe lida com isso? \perfil{A B C}
  \item Clientes pedem cardápio impresso? Com que frequência? Como vocês atendem a lei municipal que exige o cardápio físico? \perfil{A B C}
  \item Vocês usam algum dado gerado pelo sistema (itens mais vendidos, horários de pico) para tomar decisões? \perfil{A C}
\end{enumerate}

%-------------------------------------------------------------
\bloco{Bloco 6, Futuro}{todos}{5 min}
\begin{enumerate}[resume=perguntas]
  \item Daqui a cinco anos, como você imagina o atendimento na sua casa? \perfil{A B C}
  \item Se pudesse mudar uma única coisa nos sistemas de cardápio digital que existem hoje, o que seria? \perfil{A B C}
\end{enumerate}

%-------------------------------------------------------------
\bloco{Bloco 7, Encerramento}{todos}{2 min}
\begin{enumerate}[resume=perguntas]
  \item Tem algo importante sobre esse assunto que eu não perguntei? \perfil{A B C}
  \item Posso divulgar o questionário dos clientes aqui (QR Code no balcão ou nas redes da casa)? \perfil{A B C}
\end{enumerate}

\begin{tcolorbox}[colback=gray!6,colframe=gray!50,title=Logo depois da entrevista,fonttitle=\bfseries]
\caixa Agradecer e combinar o envio do resumo dos resultados.\\
\caixa Fotografar a interface ou o equipamento (se autorizado), para a análise de usabilidade.\\
\caixa Anotar em 5 minutos: impressões gerais, o que mais surpreendeu, contradições e o que não ficou gravado.\\
\caixa Salvar o áudio com o nome \texttt{perfil-X\_estabelecimento\_AAAA-MM-DD}.
\end{tcolorbox}

\newpage
\textbf{Anotações}\par
\vspace{0.5em}
\newcount\linhas \linhas=0
\loop\ifnum\linhas<24 \rule{\linewidth}{0.3pt}\par\vspace{0.7em}\advance\linhas by 1 \repeat
\end{document}
